\chapter{Proof Complexity and SAT Solvers}
\label{ch:proofcomplexity}

The practical way to solve $\NP$ complete problems is not Monte Carlo but search: the DPLL procedure, conflict driven clause learning, and the industrial SAT solvers that routinely handle instances with millions of variables. This chapter describes how they work, what proof complexity says about their limits, and why they and the Monte Carlo methods of this book fail on the same instances for related reasons. It also records what the queens look like to a SAT solver, which is a useful calibration.

\section{Resolution and DPLL}

A SAT solver decides whether a CNF formula is satisfiable. If it is, the solver outputs an assignment, which is a certificate. If it is not, the solver's trace is a proof of unsatisfiability, and the proof is in a specific proof system, resolution: from clauses $C \vee x$ and $D \vee \neg x$ derive $C \vee D$, and the formula is unsatisfiable if the empty clause is derivable. Resolution is complete for propositional unsatisfiability and is the system in which most practical solvers reason.

The Davis Putnam Logemann Loveland procedure of 1962 searches assignments by branching on a variable, propagating the unit clauses that the choice forces, and backtracking on conflict. Its search tree is a tree resolution proof when the formula is unsatisfiable. Conflict driven clause learning, due to Marques Silva and Sakallah in 1996 and refined by Moskewicz, Madigan, Zhao, Zhang and Malik in Chaff, adds to this the derivation of a new clause at each conflict, which summarises the reason for the conflict and prevents the search from revisiting it, together with nonchronological backjumping and restarts. The learned clauses make the proof a general resolution proof rather than a tree one, and Pipatsrisawat and Darwiche showed in 2011 that CDCL with suitable decisions and restarts polynomially simulates general resolution.

\section{Lower bounds for resolution}

Haken proved in 1985 that the pigeonhole principle, the CNF encoding of the statement that $n + 1$ pigeons do not fit in $n$ holes, requires resolution proofs of size $2^{\Omega(n)}$. The result was the first superpolynomial lower bound for general resolution on a natural formula, and it says that every DPLL or CDCL solver takes exponential time on the pigeonhole formulas. Ben Sasson and Wigderson\index{Wigderson, Avi} gave in 2001 a general method, relating proof size to proof width by the inequality that size is at least $2^{\Omega((w - w_0)^2 / n)}$ where $w$ is the minimum width of a refutation and $w_0$ the width of the formula, and used it to recover Haken's bound and to prove that random $k$ CNF formulas at clause densities above the satisfiability threshold require exponential resolution proofs with high probability, a result first obtained by Chv\'atal and Szemer\'edi in 1988 for sufficiently high densities.

The consequence for solvers is sharp. On random $3$ SAT at density above about $4.27$, which by the cavity method prediction is unsatisfiable with high probability, every CDCL solver takes exponential time to prove it. On random $3$ SAT below the threshold, which by the same prediction is satisfiable, the solvers find an assignment quickly at low density and struggle near the threshold, where, by the cavity method predictions that Achlioptas and Coja Oghlan proved for large $k$, the solution space shatters into clusters, and the struggle is the same phenomenon that stops the Monte Carlo chains: the local structure of the search space has no gradient towards the solutions.

\section{Stronger proof systems}

Resolution is weak. Cutting planes, which reason with linear inequalities over the integers, prove the pigeonhole principle in polynomial size. Extended resolution, which allows the introduction of new variables defined by formulas, is polynomially equivalent to extended Frege, the strongest known system, and no superpolynomial lower bound is known for it. Frege systems, the textbook propositional calculi, are in between, and lower bounds for them are open. The question whether there is a propositional proof system in which every tautology has a polynomial size proof is equivalent, by Cook\index{Cook and Reckhow theorem} and Reckhow's theorem of 1979, to $\NP = \mathrm{co}\NP$; so proof complexity is the study of $\NP$ versus $\mathrm{co}\NP$ one proof system at a time, and it has reached, in more than forty years, resolution, cutting planes, polynomial calculus, bounded depth Frege and a few relatives.

\section{The queens as a SAT instance}

The $N$ queens problem encodes as a CNF formula on $N^2$ variables $q_{ij}$, one per square, with clauses saying each row has at least one queen, and no two queens share a row, column or diagonal; the at most one constraints are pairwise and give $O(N^3)$ clauses. The unconditioned problem is trivially satisfiable and one expects a solver to find a solution quickly for $N$ in the hundreds (we have not run this experiment), since the constraints propagate well. The completion problem on random partial boards is also expected to be easy for solvers at most densities, and hard only on the reduction instances, where it is hard because the formula being encoded is hard.

Counting is different. A $\sharp\mathrm{SAT}$ solver, such as those based on component caching and knowledge compilation, counts satisfying assignments by decomposing the formula into independent components and multiplying their counts, with caching of subformulas. On the queens the components do not decompose, since every queen constrains every other along some line, and the counters fall back to enumeration. The exact counts to $N = 27$ were obtained by specialised enumeration with symmetry reduction and hardware acceleration, not by general counters, and we expect the general counters to give up at modest $N$ (we have not benchmarked them).

\section{Search versus sampling}

A SAT solver searches for one solution; a Monte Carlo chain samples among many. On an unsatisfiable or barely satisfiable instance the two problems are the same and both methods fail for the same reason, the absence of local structure that leads towards a solution. On a richly satisfiable instance the two problems diverge: the solver stops at the first solution and the chain must visit all of them in proportion. For counting the chain's task is the right one, since one wants the number of solutions and not any one of them.

Proof complexity gives one thing to Monte Carlo that it does not have on its own: a theory of lower bounds for a specific algorithmic method on specific instances. Haken's theorem says what DPLL cannot do. There is no analogous theorem that says what the swap chain cannot do on the queens; the torpid mixing result of Gore and Jerrum\index{Gore and Jerrum} is of this kind for another chain and another model, and the open problem of Chapter~\ref{ch:open} on the bottom rung for conditioned boards asks for one here. Such a theorem would be the Monte Carlo analogue of Haken: not $\Pc \neq \NP$, but a proof that the natural method fails, with the reason named.

\section{Hybrid methods}

The most effective practical methods combine search and sampling. Survey propagation, due to M\'ezard, Parisi and Zecchina in 2002, uses a message passing scheme derived from the cavity method of statistical physics to estimate, for each variable, the probability that it is frozen to a value in the clusters of solutions, fixes the most polarised variables, and recurses; it solves random $3$ SAT closer to the threshold than any other method discussed here. Its success is the clearest demonstration that the cluster structure of the solution space, which is what stops both DPLL and the Markov chains, can be exploited rather than fought if one has a model of it. For the queens no cavity method has been developed; the interactions are not sparse, since every queen interacts with every other, and the cavity method needs a locally treelike factor graph. Whether a dense analogue exists, along the lines of the replica method for fully connected spin glasses, is a question we have not pursued and which may be the right route to the structure of the queens level sets.
